\documentclass[12pt,a4paper]{article}
\usepackage{amsmath,amssymb,amsthm}
\usepackage{geometry}
\usepackage{enumitem}
\usepackage{xcolor}
\usepackage{mdframed}
\usepackage{multicol}
\usepackage{array}
\geometry{margin=2cm}

\definecolor{anscolor}{RGB}{0,120,0}
\definecolor{partbg}{RGB}{230,240,255}
\newcommand{\ans}[1]{\quad\textcolor{anscolor}{\textbf{[Answer: #1]}}}
\newcommand{\shead}[1]{\vspace{1em}{\large\bfseries\color{blue}#1}\vspace{0.3em}\hrule\vspace{0.5em}}

\begin{document}

\begin{center}
{\LARGE\bfseries CSIR-UGC NET December 2025}\\[4pt]
{\large Mathematical Sciences (Code 704) \quad 18 December 2025, Shift 2}\\[4pt]
{\large Complete Question Paper with Final Answer Key}
\end{center}
\hrule\vspace{1em}

%% ============================================================
\shead{PART A --- General Aptitude (Q.1--Q.20)}
%% ============================================================
\begin{mdframed}[backgroundcolor=partbg,linewidth=0pt]
Marking: $+2$ correct, $-0.5$ wrong. Attempt any 15 of 20.
\end{mdframed}
\vspace{0.5em}

\textbf{Q.1} Five students graduated from a college, not all in the same year, after each has studied for four years. If batchmates Jiten and Anwar were between Ramesh and Prakash but senior to Sam while Ramesh had left the college before Jiten took admission, then it is certain that
\begin{enumerate}[label=\arabic*.]
\item Anwar was the most senior \item Ramesh was the most senior
\item Sam was the most junior \item Prakash was the most junior
\end{enumerate}\ans{2}

\textbf{Q.2} Which among the following cities can be said most appropriately to bear the same relation to \textit{Tamil Nadu} that \textbf{Pune} bears to \textit{Maharashtra}; \textbf{Surat} to \textit{Gujarat} and \textbf{Asansol} to \textit{West Bengal}?
\begin{enumerate}[label=\arabic*.]
\item Tirupati \item Mysore \item Chennai \item Coimbatore
\end{enumerate}\ans{4}

\textbf{Q.3} The geometric mean of 100 observations is 25. If each observation is multiplied by 4, what will be the new geometric mean?
\begin{enumerate}[label=\arabic*.]
\item 100 \item 50 \item 25 \item $(25\times4)^{1/2}$
\end{enumerate}\ans{1}

\textbf{Q.4} Alloy A is formed by mixing iron (Fe) and nickel (Ni) in the ratio 3:4, while alloy B is formed by mixing Fe and Ni in the ratio 9:5. If equal quantities of alloys A and B are melted together to form a new alloy C, what will be the ratio of Fe to Ni in the alloy C?
\begin{enumerate}[label=\arabic*.]
\item 4:3 \item 5:3 \item 15:13 \item 13:9
\end{enumerate}\ans{3}

\textbf{Q.5} Two identical metal bars are heated to different temperatures and allowed to cool in the same surroundings. Which one of the following figures correctly shows their cooling curves? (Options A, B, C, D --- curves showing temperature vs.\ time for two bars starting at different initial temperatures.)
\begin{enumerate}[label=\arabic*.]
\item A \item B \item C \item D
\end{enumerate}\ans{3}

\textbf{Q.6} The following figure shows densities and masses of five objects (A to E). The object with the largest volume is \underline{\hspace{1cm}}.
\begin{enumerate}[label=\arabic*.]
\item A \item B \item D \item E
\end{enumerate}\ans{3}

\textbf{Q.7} In an exam, questions of three difficulty levels hard, medium, and easy fetch respectively 7, 3, and 2 marks if correct and 0 if incorrect. Three students got 30 marks each but in three different ways, though the total number of questions correctly answered by each student was the same. Then what could be the total number of questions correctly answered by each of these students?
\begin{enumerate}[label=\arabic*.]
\item 12 \item 10 \item 9 \item 6
\end{enumerate}\ans{2}

\textbf{Q.8} The following plot shows temperature as a function of $x$ and $y$. Along which path is the temperature change minimum?
\begin{enumerate}[label=\arabic*.]
\item $x=\text{constant}$ or $y=\text{constant}$
\item $\dfrac{y}{x^2}=\text{constant}$
\item $y^2+x^2=\text{constant}$
\item $yx=\text{constant}$
\end{enumerate}\ans{3}

\textbf{Q.9} The minimum height of a plane vertical mirror that will allow a 6-feet tall person to see himself fully in it
\begin{enumerate}[label=\arabic*.]
\item depends on the distance between the person and the mirror
\item is 3 feet \item is 4.5 feet \item is 6 feet
\end{enumerate}\ans{2}

\textbf{Q.10} Three periodic events repeat every 24 seconds, 54 seconds, and 56 seconds. If they coincide at 10:20:00, when will they next coincide?
\begin{enumerate}[label=\arabic*.]
\item 10:35:12 \item 10:45:20 \item 10:45:12 \item 10:35:20
\end{enumerate}\ans{3}

\textbf{Q.11} The numbers (in millions) of coconut, guava, mango and apple trees in a region in 2010 and 2020 are shown in the figure. The maximum relative change in numbers was for
\begin{enumerate}[label=\arabic*.]
\item coconut trees \item guava trees \item mango trees \item apple trees
\end{enumerate}\ans{1}

\textbf{Q.12} Some, but not all, faces of a six-faced cubical fair die are painted red (R) and the remaining green (G); and the die is thrown until red faces come up on top 4 times. Consider the following sequences of colours listed left to right as they appear on the top.\\
A: \textbf{GRRRR}\quad B: \textbf{GRGRRR}\\
Which one of the following is true?
\begin{enumerate}[label=\arabic*.]
\item A is more probable than B
\item B is more probable than A
\item Both have the same probability
\item Whether A or B is more probable depends upon how many faces are painted green
\end{enumerate}\ans{1}

\textbf{Q.13} (Pouring problem) From a 12-litre full vessel, transfer 6 litres to an 8-litre empty vessel, given that a 5-litre empty vessel is also available. The minimum number of pourings required is:
\begin{enumerate}[label=\arabic*.]
\item 4 \item 5 \item 6 \item 7
\end{enumerate}\ans{3}

\textbf{Q.14} A lady bought some apples, each costing Rs.\,25, and some bananas each costing Rs.\,6, for a total of Rs.\,378. In how many ways could she have chosen the numbers of apples and bananas?
\begin{enumerate}[label=\arabic*.]
\item 1 \item 2 \item 3 \item 4
\end{enumerate}\ans{2}

\textbf{Q.15} Suppose $a_1,a_2,\ldots,a_{300}$ are integers such that $a_{i-1}+a_i+a_{i+1}=2025$ for all $i=2,3,\ldots,299$. If $a_7=-5$, $a_9=37$, then the value of $a_{106}$ is
\begin{enumerate}[label=\arabic*.]
\item 1993 \item 37 \item $-5$ \item 2030
\end{enumerate}\ans{3}

\textbf{Q.16} How many 5-digit numbers can be formed from the digits 0, 2, 3, 4, 6, 7 and 9, using each at most once, which are divisible by 5?
\begin{enumerate}[label=\arabic*.]
\item 120 \item 240 \item 360 \item 720
\end{enumerate}\ans{3}

\textbf{Q.17} A recent survey suggests that the total fertility rate in a country has fallen below 2.1, the population replacement ratio. This necessarily implies that the
\begin{enumerate}[label=\arabic*.]
\item infant mortality rate has increased reducing the net fertility ratio.
\item total population will decline.
\item population of young people is going to increase with a faster rate in the long run if the same status continues.
\item proportion of elderly people is going to decrease in the long run if the same status continues.
\end{enumerate}\ans{2}

\textbf{Q.18} In a class, 40\% and 20\% students passed in Mathematics and Physics, respectively, and 10\% students passed in both subjects. What is the probability of a randomly selected student to have passed in Physics if the student already passed in Mathematics?
\begin{enumerate}[label=\arabic*.]
\item $1/2$ \item $1/20$ \item $1/4$ \item $2/25$
\end{enumerate}\ans{3}

\textbf{Q.19} The value of
\[1+\!\left(\frac{1}{2^1}+\frac{1}{3}\right)\!+\!\left(\frac{1}{2^2}+\frac{1}{5}+\frac{1}{6}+\frac{1}{7}\right)\!+\cdots+\!\left(\frac{1}{2^9}+\cdots+\frac{1}{1023}\right)\]
lies between
\begin{enumerate}[label=\arabic*.]
\item 2 and 10 \item 11 and 20 \item 21 and 30 \item 31 and 40
\end{enumerate}\ans{1}

\textbf{Q.20} In a community, some artists are teachers, no teacher is a painter, all painters are artists, and all teachers are professionals. Then it can be definitely asserted that
\begin{enumerate}[label=\arabic*.]
\item no painter is a professional \item all artists are professionals
\item no professionals are teachers \item some artists are professionals
\end{enumerate}\ans{4}

%% ============================================================
\shead{PART B --- Mathematics MCQ (Q.21--Q.60)}
%% ============================================================
\begin{mdframed}[backgroundcolor=partbg,linewidth=0pt]
Marking: $+3$ correct, $-1$ wrong. Attempt any 25 of 40.
\end{mdframed}
\vspace{0.5em}

\textbf{Q.21} Let $X$ be a Binomial$(n,p)$ random variable, where $n\in\{5,6\}$ and $p\in\{\frac{1}{4},\frac{3}{4}\}$. If $X=3$ is observed, then the maximum likelihood estimate of $(n,p)$ is
\begin{enumerate}[label=\arabic*.]
\item $\bigl(5,\tfrac{1}{4}\bigr)$ \item $\bigl(5,\tfrac{3}{4}\bigr)$ \item $\bigl(6,\tfrac{3}{4}\bigr)$ \item $\bigl(6,\tfrac{1}{4}\bigr)$
\end{enumerate}\ans{2}

\textbf{Q.22} Let $X=\{1,\ldots,17\}$ and $S_{17}$ be the group of permutations of $X$. For a subgroup $G$ of $S_{17}$ and $x\in X$, let $S_G(x)=\{\sigma\in G\mid\sigma(x)=x\}$. Which of the following statements is true for every subgroup $G$ of $S_{17}$?
\begin{enumerate}[label=\arabic*.]
\item The number of pairs $(\sigma,x)\in G\times X$ such that $\sigma(x)=x$ is strictly greater than $\displaystyle\sum_{x\in X}|S_G(x)|$.
\item The number $\dfrac{1}{|G|}\displaystyle\sum_{x\in X}|S_G(x)|$ is always an integer.
\item For all $x\in X$, $S_G(x)$ is a normal subgroup of $G$.
\item For all $x,y\in X$, $S_G(x)$ is isomorphic to $S_G(y)$.
\end{enumerate}\ans{2}

\textbf{Q.23} Let $I$ denote the $3\times3$ identity matrix. Given any three distinct matrices $A,B,C\in M_3(\mathbb{R})$, which of the following statements is necessarily true?
\begin{enumerate}[label=\arabic*.]
\item There exists $D\in M_3(\mathbb{R})$ and $f\in\mathbb{R}[X]$ satisfying $f(A)=f(B)=f(C)=0$ and $f(D)=I$.
\item There exists $D\in M_3(\mathbb{R})$ and $f\in\mathbb{R}[X]$ satisfying $f(A)=f(B)=0$, $f(C)=I$ and $f(D)=I$.
\item There exists $D\in M_3(\mathbb{R})$ and $f\in\mathbb{R}[X]$ satisfying $f(A)=0$, $f(B)=f(C)=I$ and $f(D)\neq0$.
\item There exists $D\in M_3(\mathbb{R})$ and $f\in\mathbb{R}[X]$ satisfying $f(A)=f(B)=0$, $f(C)=I$ and $f(D)=0$.
\end{enumerate}\ans{1}

\textbf{Q.24} Let $X$ and $Y$ be topological spaces. Consider the statement \textbf{S}: For every open subset $U\subseteq X\times Y$ and every $x\in X$ such that $\{x\}\times Y\subseteq U$, there is a neighbourhood $W$ of $x$ in $X$ such that $W\times Y\subseteq U$. Which of the following statements is true?
\begin{enumerate}[label=\arabic*.]
\item If $X\times Y$ is Hausdorff, then \textbf{S} is true.
\item If $Y$ is connected, then \textbf{S} is true.
\item If $X\times Y$ is regular, then \textbf{S} is true.
\item If $Y$ is compact, then \textbf{S} is true.
\end{enumerate}\ans{4}

\textbf{Q.25} Let $M,N\in M_3(\mathbb{C})$ be such that $\operatorname{Trace}(M^k)=\operatorname{Trace}(N^k)$, $1\le k\le3$. Which of the following statements is necessarily true?
\begin{enumerate}[label=\arabic*.]
\item $M^2=N^2$ \item $M^3=N^3$
\item The characteristic polynomials of $M$ and $N$ are the same.
\item The minimal polynomials of $M$ and $N$ are the same.
\end{enumerate}\ans{3}

\textbf{Q.26} Let $\{X_n:n\ge0\}$ be any homogeneous Markov Chain on the state space $S=\{1,2,3,4\}$ having the transition probability matrix
\[P=\begin{pmatrix}\tfrac{1}{4}&0&\tfrac{2}{3}&\tfrac{1}{12}\\0&1&0&0\\\tfrac{1}{12}&0&\tfrac{1}{4}&\tfrac{2}{3}\\\tfrac{2}{3}&0&\tfrac{1}{12}&\tfrac{1}{4}\end{pmatrix}.\]
Which of the following statements about stationary distributions is true?
\begin{enumerate}[label=\arabic*.]
\item There is no stationary distribution \item Stationary distribution exists and is unique
\item There are exactly two stationary distributions \item There are infinitely many stationary distributions
\end{enumerate}\ans{4}

\textbf{Q.27} Let $\mathbb{F}_q$ be a finite field with $q$ elements. For $n\ge2$, let $A$ be a $2n\times2n$ matrix with entries in $\mathbb{F}_q$ such that $\operatorname{rank}(A)=n$. Let $W=\{v\in\mathbb{F}_q^{2n}\mid Av=0\}$. Which of the following is necessarily the number of $(n+2)$-dimensional subspaces of $\mathbb{F}_q^{2n}$ that contain $W$?
\begin{enumerate}[label=\arabic*.]
\item $1$
\item $\dfrac{(q^{2n}-q^n)(q^{2n}-q^{n+1})}{q^n}$
\item $\dfrac{(q^{2n}-1)(q^{2n}-q)\cdots(q^{2n}-q^{n-1})}{q^n}$
\item $\dfrac{(q^n-1)(q^n-q)}{(q^2-1)(q^2-q)}$
\end{enumerate}\ans{4}

\textbf{Q.28} Let $\mathbb{D}=\{z\in\mathbb{C}:|z|<1\}$ and $X=\bigl\{f:\mathbb{D}\to\mathbb{C}\mid f\text{ is holomorphic and }f(2z)=\dfrac{f(z)}{1-(f(z))^2}\text{ for all }|z|<\tfrac{1}{2}\bigr\}$. Which of the following is true?
\begin{enumerate}[label=\arabic*.]
\item $X$ is uncountable. \item $X$ is infinite and countable.
\item Every element of $X$ is an open map. \item $X$ has exactly one element.
\end{enumerate}\ans{4}

\textbf{Q.29} Let $u(x,t)$ be the solution of $u_t=16u_{xx}+2$, $0<x<7$, $t>0$, satisfying $u_x(0,t)=u(7,t)=0$, $t>0$, and $u(x,0)=0$, $0<x<7$. Then which of the following is true?
\begin{enumerate}[label=\arabic*.]
\item For each $x\in(0,7)$, $u(x,t)\to0$ as $t\to\infty$.
\item For each $x\in(0,7)$, $u(x,t)\to\dfrac{x^2(7-x)^2}{16}$ as $t\to\infty$.
\item For each $x\in(0,7)$, $u(x,t)\to\dfrac{x^2}{16}(7-x)$ as $t\to\infty$.
\item For each $x\in(0,7)$, $u(x,t)\to\dfrac{1}{16}(49-x^2)$ as $t\to\infty$.
\end{enumerate}\ans{4}

\textbf{Q.30} For a finite group $G$, let $S(G)$ denote the number of subgroups of $G$. Which of the following statements is necessarily true?
\begin{enumerate}[label=\arabic*.]
\item Let $G$ and $G'$ be finite groups such that $S(G)=S(G')$. Then $G$ is isomorphic to $G'$.
\item If $S(G)=4$, then $|G|=p^m$ for some prime $p$ and positive integer $m$.
\item If $S(G)=5$, then $G$ is a cyclic group.
\item For every positive integer $n$, there exists a finite group $G$ such that $S(G)=n$.
\end{enumerate}\ans{4}

\textbf{Q.31} Let $X$, $Y$, and $Z$ be independent Normal random variables with means $-1$, $0$, and $1$, respectively, and variances $1$, $1$, and $3$, respectively. Which of the following random variables has a Cauchy distribution with location parameter 0 and scale parameter 1?
\begin{enumerate}[label=\arabic*.]
\item $\dfrac{X-1}{|Y|}$ \item $\dfrac{X+2Y+Z}{X-2Y+Z}$ \item $\dfrac{Z-1}{Y}$ \item $\dfrac{X-Y}{X+Y}$
\end{enumerate}\ans{2}

\textbf{Q.32} Suppose $y(x)$ is the extremal of the variational problem $J(y)=\int_0^4\dfrac{(y')^2}{y^2}\,dx$ subject to $y(0)=1$, $y(4)=e^8$. Then which of the following is true?
\begin{enumerate}[label=\arabic*.]
\item $y(\log_e2)=2$ \item $y(\log_e3)=9$ \item $y(\log_e4)=4$ \item $y(\log_e5)=5$
\end{enumerate}\ans{2}

\textbf{Q.33} For a non-negative real number $a$, let $\sqrt{a}$ denote its non-negative square-root. Consider $f:\mathbb{R}\to\mathbb{R}$ given by $f(x)=x\!\left(\sqrt{x^2+3}-\sqrt{x^2+2}\right)$. Which of the following is true?
\begin{enumerate}[label=\arabic*.]
\item $\lim_{x\to\infty}f(x)=\infty$ \item $\lim_{x\to\infty}f(x)=0$
\item $\lim_{x\to\infty}f(x)=1$ \item $\lim_{x\to\infty}f(x)=\dfrac{1}{2}$
\end{enumerate}\ans{4}

\textbf{Q.34} Consider the following statements:\\
\textbf{(P)} The initial value problem $y'=f(y)$, where $f(y)=\begin{cases}y\cos\frac{1}{y}, & y\neq0\\0, & y=0\end{cases}$, $y(0)=0$, has at most one solution.\\
\textbf{(Q)} If $z(x)$ is the solution of $z'=\dfrac{(1-z^4)^{100}}{1+z^2}$, $z(0)=0$, then for each $\alpha\in(0,2)$, there exists $x_\alpha\in\mathbb{R}$ such that $z(x_\alpha)=1-\alpha$.\\
Then which of the following is true?
\begin{enumerate}[label=\arabic*.]
\item Both (P) and (Q) are true. \item (P) is true, (Q) is FALSE.
\item (P) is FALSE, (Q) is true. \item Both (P) and (Q) are FALSE.
\end{enumerate}\ans{1}

\textbf{Q.35} Consider the IBVP: $u_t+u_x=0$, $0<x<\infty$, $t>0$; $u(x,0)=e^x$, $0<x<\infty$; $u(0,t)=1-\sin t$, $t>0$. Then which of the following is true?
\begin{enumerate}[label=\arabic*.]
\item There exists a unique solution $u(x,t)$ such that $u(1,t)=e-\sin t$, for all $t<1$.
\item There does NOT exist a solution $u(x,t)$ such that $u(1,t)=e-\sin t$, for all $t>1$.
\item There exists a solution $u(x,t)$ such that $u(x,t)=e^{x-t}$, for all $t>x$.
\item There exists a solution $u(x,t)$ such that $u(x,t)=e-\sin(t-x)$, for all $t<x$.
\end{enumerate}\ans{2}

\textbf{Q.36} For a non-negative real number $a$, let $\sqrt{a}$ denote its non-negative square-root. Consider the sequence $\{x_n\}_{n\ge1}$ defined by $x_1=1$ and $x_{n+1}=\sqrt{1+\frac{x_n^2}{n}}$ for $n\ge1$. Which of the following is true?
\begin{enumerate}[label=\arabic*.]
\item $\limsup_{n\to\infty}x_n=1$ \item $\limsup_{n\to\infty}x_n=\sqrt{2}$
\item $\limsup_{n\to\infty}x_n=2$ \item $\limsup_{n\to\infty}x_n=1+\sqrt{2}$
\end{enumerate}\ans{1}

\textbf{Q.37} Consider the linear programming problem: Maximize $3x_1+4x_2$ subject to $3x_1+2x_2\le12$, $3x_1+5x_2\le15$, $2x_1-x_2\ge0$, $x_2\le2$, $x_1,x_2\ge0$. Which of the following values is the optimum value of the objective function in the feasible region?
\begin{enumerate}[label=\arabic*.]
\item 11 \item 13 \item 14 \item 15
\end{enumerate}\ans{3}

\textbf{Q.38} For $a\in\mathbb{R}$, let $y_1(x)$ and $y_2(x)$ be solutions of $y''+(e^{x^2}+\cos x)y=0$ such that $y_1(0)=3$, $y_1'(0)=-1$, $y_2(0)=-5$, $y_2'(0)=a$. If $W(y_1,y_2)(\tfrac{1}{2})=4$, then the value of $a$ is
\begin{enumerate}[label=\arabic*.]
\item 2 \item 3 \item 4 \item $-3$
\end{enumerate}\ans{2}

\textbf{Q.39} Let $\alpha,\beta\in(0,\infty)$. Consider the infinite series $\displaystyle\sum_{n\ge4}\dfrac{1}{n(\log_e n)^\alpha(\log_e(\log_e n))^\beta}$. Which of the following is true?
\begin{enumerate}[label=\arabic*.]
\item The series converges for all $\alpha\in(0,1)$ and for all $\beta\in(1,\infty)$.
\item The series converges for $\alpha=1$ and $\beta=1$.
\item The series converges for $\alpha=1$ and for all $\beta\in(1,\infty)$.
\item The series converges for all $\alpha\in(0,\infty)$ and for all $\beta\in(0,\infty)$.
\end{enumerate}\ans{3}

\textbf{Q.40} Let $X=(X_1,X_2,X_3)^T$ be a $3\times1$ random vector with $E(X)=(3,2,1)^T$ and $\operatorname{Cov}(X)=\Sigma=\begin{pmatrix}3&-2&0\\-2&3&-2\\0&-2&3\end{pmatrix}$. Suppose $Y=(Y_1,Y_2,Y_3)^T=\tfrac{1}{\sqrt{3}}X$. Then the multiple correlation coefficient between $Y_1$ and $(Y_2,Y_3)$ equals
\begin{enumerate}[label=\arabic*.]
\item $\dfrac{2}{\sqrt{5}}$ \item $\sqrt{\dfrac{2}{5}}$ \item $\dfrac{2}{\sqrt{3}}$ \item $\dfrac{1}{\sqrt{3}}$
\end{enumerate}\ans{1}

\textbf{Q.41} Let $S$ be a finite subset of $\mathbb{C}$ containing 0. Let $f:\mathbb{C}\setminus S\to\mathbb{C}$ be holomorphic with a simple pole at 0. For $R>0$, let $\gamma_R(t)=Re^{2\pi it}$, $t\in[0,1]$. Which of the following is true?
\begin{enumerate}[label=\arabic*.]
\item The function $f(1/z)$ has a removable singularity at 0.
\item The function $f(1/z)$ has an essential singularity at 0.
\item If $\int_{\gamma_R}f(z)\,dz=0$ for some $R>0$, then $S$ is not a singleton set.
\item If $S$ is not a singleton set, then $\int_{\gamma_R}f(z)\,dz=0$ for some $R>0$ such that $S\cap\{z\in\mathbb{C}:|z|=R\}=\emptyset$.
\end{enumerate}\ans{3}

\textbf{Q.42} Let $h$ be holomorphic on $\mathbb{C}\setminus\{0\}$ such that $\lim_{|z|\to\infty}h(z)=0$. For every $n\ge1$, consider $f_n(z)=\sum_{k=1}^n h(z^k)$. Let $\mathbb{D}=\{z\in\mathbb{C}:|z|<1\}$. Which of the following is necessarily true?
\begin{enumerate}[label=\arabic*.]
\item For all $z$ with $|z|\ge1$, the sequence $\{f_n(z)\}_{n\ge1}$ converges.
\item For all $z\in\mathbb{D}\setminus\{0\}$, the sequence $\{f_n(z)\}_{n\ge1}$ converges.
\item The sequence $\{f_n\}_{n\ge1}$ converges pointwise to a holomorphic function on $\mathbb{D}\setminus[0,1)$.
\item The sequence $\{f_n\}_{n\ge1}$ converges pointwise to a holomorphic function on $\{z\in\mathbb{C}:|z|>1\}$.
\end{enumerate}\ans{4}

\textbf{Q.43} Consider a finite population of size $N=100$. Let $T_1$ be the sample mean under SRSWR of size $n$ ($1<n<N$). Let $T_2$ be the sample mean under SRSWOR of size $n$. If $\operatorname{Var}(T_1)=9\operatorname{Var}(T_2)$, then $n$ equals
\begin{enumerate}[label=\arabic*.]
\item 33 \item 69 \item 89 \item 93
\end{enumerate}\ans{3}

\textbf{Q.44} Let $\gamma_R(t)=2+i+Re^{2\pi it}$ for $t\in[0,1]$ and $R=1,2$. Which of the following is true?
\begin{enumerate}[label=\arabic*.]
\item $\displaystyle\int_{\gamma_1}\tan(z)\,dz=2\pi i$ \item $\displaystyle\int_{\gamma_1}\tan(z)\,dz=-2\pi i$
\item $\displaystyle\int_{\gamma_2}\tan(z)\,dz=2\pi i$ \item $\displaystyle\int_{\gamma_2}\tan(z)\,dz=-2\pi i$
\end{enumerate}\ans{4}

\textbf{Q.45} Let $B$ be a $4\times4$ positive-definite real symmetric matrix which is not the identity matrix. Consider the inner product on $\mathbb{R}^4$ given by $\langle v,w\rangle=v^\top Bw$. Which of the following is \textbf{FALSE}?
\begin{enumerate}[label=\arabic*.]
\item If $v$ and $w$ are eigenvectors of $B$ for distinct eigenvalues, then $\langle v,w\rangle=0$.
\item If $v$ is an eigenvector of $B$ and $w\neq0$ is such that $\langle w,v\rangle=0$, then $w$ is an eigenvector of $B$.
\item For every subspace $W\subseteq\mathbb{R}^4$, $W+W^\perp=\mathbb{R}^4$.
\item If $W\subseteq\mathbb{R}^4$ is an eigenspace of $B$, then $W$ has an orthonormal basis.
\end{enumerate}\ans{2}

\textbf{Q.46} Let $u(x,t)$ be the solution of $u_{tt}-u_{xx}=e^{-t}$, $x\in\mathbb{R}$, $t>0$, $u(x,0)=\cos x$, $u_t(x,0)=0$, $x\in\mathbb{R}$. Then which of the following is true?
\begin{enumerate}[label=\arabic*.]
\item For each $x_0\in\mathbb{R}$, $e^tu(x_0,t)\to0$ as $t\to\infty$.
\item For each $x_0\in\mathbb{R}$, $|u(x_0,t)|\to0$ as $t\to\infty$.
\item For each $x_0\in\mathbb{R}$, there exists $t_0>0$ such that $u(x_0,t_0)\ge4$.
\item For each $x_0\in\mathbb{R}$, there exists $t_0>0$ such that $u(x_0,t_0)\le-4$.
\end{enumerate}\ans{3}

\textbf{Q.47} Let $f:(0,\infty)\to(0,\infty)$ be uniformly continuous. Define $g:(0,\infty)\to(0,\infty)$ by $g(x)=\begin{cases}x^2&\text{if }0<x\le1\\x&\text{if }x>1\end{cases}$. Which of the following is true?
\begin{enumerate}[label=\arabic*.]
\item $f\circ g$ is uniformly continuous, but $g\circ f$ is not uniformly continuous.
\item $g\circ f$ is uniformly continuous, but $f\circ g$ is not uniformly continuous.
\item Neither $f\circ g$ nor $g\circ f$ is uniformly continuous.
\item Both $f\circ g$ and $g\circ f$ are uniformly continuous.
\end{enumerate}\ans{4}

\textbf{Q.48} If $p(x)$ is the interpolating polynomial for the data
\[\begin{array}{c|ccccc}x&-2&-1&0&1&2\\\hline y&-1&3&1&-1&3\end{array}\]
then the value of $p\!\left(\tfrac{1}{2}\right)$ is
\begin{enumerate}[label=\arabic*.]
\item $\dfrac{-3}{8}$ \item $\dfrac{-5}{8}$ \item $\dfrac{5}{8}$ \item $\dfrac{3}{8}$
\end{enumerate}\ans{1}

\textbf{Q.49} Let $S$ be the subset of $(0,1)$ consisting of all $\alpha\in(0,1)$ whose infinite decimal expansion $\alpha=0.a_1a_2a_3\cdots$ satisfies $a_i\in\{0,2,4\}$ for all $i\ge1$. Which of the following is true?
\begin{enumerate}[label=\arabic*.]
\item There is a bijective map from $\mathbb{N}$ to $S$.
\item There is a surjective map from $S$ onto $(0,1)$.
\item There is a bijective map from $\mathbb{N}$ to $(0,1)\setminus S$.
\item $S$ is countable and $(0,1)\setminus S$ is uncountable.
\end{enumerate}\ans{2}

\textbf{Q.50} Consider the quadratic form $f(x,y,z)=\begin{bmatrix}x&y&z\end{bmatrix}A\begin{bmatrix}x\\y\\z\end{bmatrix}$ where $A$ is an invertible $3\times3$ symmetric matrix over $\mathbb{Q}$. Assume there exists $(\alpha,\beta,\gamma)\in\mathbb{C}^3\setminus\{(0,0,0)\}$ such that $f(\alpha,\beta,\gamma)=0$. Which of the following is necessarily true?
\begin{enumerate}[label=\arabic*.]
\item There exists $(a,b,c)\in\mathbb{R}^3\setminus\{(0,0,0)\}$ such that $f(a,b,c)=0$.
\item If there exists $(a,b,c)\in\mathbb{R}^3\setminus\{(0,0,0)\}$ with $f(a,b,c)=0$, then $(a,b,c)\in\mathbb{Q}^3$.
\item $\{(a,b,c)\in\mathbb{Z}^3\mid f(a,b,c)=0\}$ is finite.
\item If there exists $(a,b,c)\in\mathbb{Q}^3\setminus\{(0,0,0)\}$ with $f(a,b,c)=0$, then $A$ has a positive eigenvalue and a negative eigenvalue.
\end{enumerate}\ans{4}

\textbf{Q.51} Let $X_1$ and $X_2$ be a random sample from Uniform$[0,\theta]$ distribution ($\theta>0$). For testing $H_0:\theta=1$ against $H_1:\theta=2$, consider a test which rejects $H_0$ if $X_1+X_2>\tfrac{4}{5}$. Then the probability of type-I error is
\begin{enumerate}[label=\arabic*.]
\item $\dfrac{8}{25}$ \item $\dfrac{13}{25}$ \item $\dfrac{17}{25}$ \item $\dfrac{22}{25}$
\end{enumerate}\ans{3}

\textbf{Q.52} Let $\{Y_n:n\ge1\}$ be i.i.d.\ with $Y_1\sim\text{Bernoulli}(\tfrac{1}{2})$. Define $Z=\sum_{n=1}^\infty\dfrac{4Y_n}{5^n}$. Then which of the following is true?
\begin{enumerate}[label=\arabic*.]
\item $P(Z\ge\tfrac{3}{5})=0.5,\ P(Z=\tfrac{4}{25})=0$
\item $P(Z\ge\tfrac{3}{5})=0.6,\ P(Z=\tfrac{4}{25})=0$
\item $P(Z\ge\tfrac{3}{5})=0.7,\ P(Z=\tfrac{4}{25})=0.16$
\item $P(Z\ge\tfrac{3}{5})=0.8,\ P(Z=\tfrac{4}{25})=0.16$
\end{enumerate}\ans{1}

\textbf{Q.53} Let $A$ be a $3\times3$ matrix over $\mathbb{C}$ with trace 1 and determinant 1. Suppose one of the eigenvalues of $A$ is 1. Which of the following is necessarily true?
\begin{enumerate}[label=\arabic*.]
\item The characteristic polynomial of $A$ has repeated roots.
\item Every eigenvalue of $A$ has absolute value 1.
\item $A$ does not have any eigenvalue on the imaginary axis.
\item $A^2$ is the identity matrix.
\end{enumerate}\ans{2}

\textbf{Q.54} Let $X|\theta\sim\text{Uniform}[0,\theta]$ and $\theta$ has an Exponential distribution with mean $\lambda$, where $\lambda>3$ is known. If the realized value of $X$ is 2025, then the posterior mode equals
\begin{enumerate}[label=\arabic*.]
\item $0$ \item $2025$ \item $\dfrac{2025}{\lambda}$ \item $2025\lambda$
\end{enumerate}\ans{2}

\textbf{Q.55} Let $X$ be a single sample from $f(x|\theta)=\begin{cases}\dfrac{2}{\theta^2}(\theta-x),&0<x<\theta\\0,&\text{otherwise}\end{cases}$ where $\theta>0$ is unknown. Which of the following is a 95\% confidence interval for $\theta$?
\begin{enumerate}[label=\arabic*.]
\item $\Bigl(X,\dfrac{X}{1-\sqrt{0.95}}\Bigr)$
\item $\Bigl(X,\dfrac{X}{0.95}\Bigr)$
\item $\Bigl(X,\dfrac{X}{\sqrt{0.95}}\Bigr)$
\item $(0.95X,X)$
\end{enumerate}\ans{1}

\textbf{Q.56} Which of the following is necessarily true?
\begin{enumerate}[label=\arabic*.]
\item The set of all finite subsets of $\mathbb{Z}$ is uncountable.
\item Let $f:\mathbb{R}\to\mathbb{R}$ be continuous and one-one. If $S\subseteq\mathbb{R}$ is countably infinite, then $f(S)$ is countably infinite.
\item Let $f:\mathbb{R}\to\mathbb{R}$ be continuous. Then there exists an uncountable subset $T\subseteq\mathbb{R}$ such that $f(T)\subseteq\mathbb{R}$ is uncountable.
\item Let $f:\mathbb{R}\to\mathbb{R}$ be continuous. If $R\subseteq\text{image}(f)$ is countable, then $f^{-1}(R)$ is countable.
\end{enumerate}\ans{2}

\textbf{Q.57} Consider the linear model $Y_1=\beta_1+\beta_2+\epsilon_1$, $Y_2=\beta_1+2\beta_2+\epsilon_2$, $Y_3=\beta_1+c\beta_2+\epsilon_3$, where $\beta_1,\beta_2\in\mathbb{R}$ are unknown, the uncorrelated errors have zero mean and variance $\sigma^2(>0)$, and $c$ is such that $\hat{\beta}_1$ and $\hat{\beta}_2$ are uncorrelated (BLUEs). Which of the following is correct for $(\operatorname{Var}(\hat{\beta}_1),\operatorname{Var}(\hat{\beta}_2))$?
\begin{enumerate}[label=\arabic*.]
\item $\Bigl(\dfrac{\sigma^2}{3},\dfrac{\sigma^2}{14}\Bigr)$ \item $(3\sigma^2,14\sigma^2)$
\item $\Bigl(\dfrac{\sigma^2}{14},\dfrac{\sigma^2}{3}\Bigr)$ \item $(14\sigma^2,3\sigma^2)$
\end{enumerate}\ans{1}

\textbf{Q.58} Suppose $u(x)$ is the solution of the integral equation $u(x)=3+\int_0^x(x-t)u(t)\,dt$. Then which of the following is true?
\begin{enumerate}[label=\arabic*.]
\item $u(\pi)=2e^\pi$ \item $u'(\pi)=e^\pi$
\item $u(\pi)+u'(\pi)=3e^\pi$ \item $u(\pi)-u'(\pi)=e^\pi$
\end{enumerate}\ans{3}

\textbf{Q.59} Let $X_1,X_2,\ldots$ be i.i.d.\ with Binomial$(3,\tfrac{1}{4})$ distribution. For $j=1,2,\ldots$, define $Y_j=\begin{cases}1,&X_j\le\sqrt{5}\\0,&\text{otherwise}\end{cases}$. If $\tfrac{1}{n}\sum_{j=1}^nY_j^2$ converges almost surely to a constant $c$ as $n\to\infty$, then $c$ equals
\begin{enumerate}[label=\arabic*.]
\item $\dfrac{1}{64}$ \item $\dfrac{37}{64}$ \item $\dfrac{63}{64}$ \item $\dfrac{27}{32}$
\end{enumerate}\ans{3}

\textbf{Q.60} For a finite group $G$, let $H_G=\{g\in G\mid g^{15}=e\}$. Which of the following is necessarily true?
\begin{enumerate}[label=\arabic*.]
\item There exists a finite group $G$ such that $|H_G|$ is even.
\item There exists a finite group $G$ such that $|H_G|=4n+1$ for some $n\ge3$.
\item For every finite group $G$, there exists a non-negative integer $n$ such that $|H_G|=4n+1$.
\item For every finite group $G$, there exists a non-negative integer $n$ such that $|H_G|=4n+3$.
\end{enumerate}\ans{2}

%% ============================================================
\shead{PART C --- Mathematics MSQ (Q.61--Q.120)}
%% ============================================================
\begin{mdframed}[backgroundcolor=partbg,linewidth=0pt]
Marking: $+4.75$ fully correct (all correct options, no wrong), $-1.25$ for any wrong option selected. Attempt any 20 of 60.
\end{mdframed}
\vspace{0.5em}

\textbf{Q.61} Let $X_1,X_2,\ldots,X_{15}$ be a random sample from Exponential distribution with pdf $f(x|\sigma)=\tfrac{1}{\sigma}e^{-x/\sigma}$, $x>0$, $\sigma>0$. Let $\bar{X}=\tfrac{1}{15}\sum X_i$. Suppose $\phi$ denotes the LRT for testing $H_0:\sigma\le1$ against $H_1:\sigma>1$ at $\alpha=0.1$. Given $\chi^2_{15,0.1}=22.307$, $\chi^2_{15,0.9}=8.547$, $\chi^2_{30,0.1}=40.256$, $\chi^2_{30,0.9}=20.599$. Then which of the following are true?
\begin{enumerate}[label=\arabic*.]
\item If $\bar{X}=0.6$, then $\phi$ does not reject $H_0$
\item If $\bar{X}=1.6$, then $\phi$ rejects $H_0$
\item If $\bar{X}=1.3$, then $\phi$ rejects $H_0$
\item If $\bar{X}=1.2$, then $\phi$ does not reject $H_0$
\end{enumerate}\ans{1, 2, 4}

\textbf{Q.62} Which of the following are true?
\begin{enumerate}[label=\arabic*.]
\item $x^2+xy^2-x-y+1$ is irreducible in $\mathbb{Q}[x,y]$.
\item $x^3+100x^2+25x+10$ is irreducible in $\mathbb{Z}[x]$.
\item $2x^2+3xy+y^2+3x+2y+1$ is irreducible in $\mathbb{Q}[x,y]$.
\item $x^4+4x^2+3$ is irreducible in $\mathbb{Z}[x]$.
\end{enumerate}\ans{1, 2}

\textbf{Q.63} Let $X_1,X_2,\ldots,X_n$ ($n\ge2$) be a random sample from Exponential distribution with pdf $f(x|\mu,\sigma)=\tfrac{1}{\sigma}e^{-(\mu-x)/\sigma}$, $x>\mu$, parameters $\mu,\sigma>0$ unknown. Let $\bar{X}_n$, $S_n^2$, $X_{1:n}$ denote sample mean, sample variance and smallest order statistic, respectively, and let $\theta=\tfrac{\sigma}{\mu}$. Then which of the following are consistent estimators of $\theta$?
\begin{enumerate}[label=\arabic*.]
\item $S_n(X_{1:n})^{-1}$
\item $(\bar{X}_n-X_{1:n})(X_{1:n})^{-1}$
\item $(\bar{X}_n-X_{1:n})(\bar{X}_n-S_n)^{-1}$
\item $S_n(\bar{X}_n)^{-1}$
\end{enumerate}\ans{1, 2, 3}

\textbf{Q.64} Consider a dataset $A=\{x_1,x_2,\ldots,x_n\}$. Create a new dataset $B=\{x_1,x_2,\ldots,x_n,-x_1,-x_2,\ldots,-x_n\}$. Which of the following are always true?
\begin{enumerate}[label=\arabic*.]
\item Mean of the observations in dataset $B$ is 0
\item Median of the observations in dataset $B$ is 0
\item Variance of the observations in dataset $B\ge$ Variance of the observations in dataset $A$
\item Range of the observations in dataset $B\ge$ Range of the observations in dataset $A$
\end{enumerate}\ans{1, 2, 3, 4}

\textbf{Q.65} Let $G$ be a group of order 8. Which of the following are necessarily true?
\begin{enumerate}[label=\arabic*.]
\item If there are no elements of order 4 in $G$, then $G$ is abelian.
\item If there are exactly two elements of order 4 in $G$, then $G$ is abelian.
\item If there are exactly six elements of order 4 in $G$, then $G$ is abelian.
\item There is an element of order 4 in $G$.
\end{enumerate}\ans{1}

\textbf{Q.66} Let $\alpha=\displaystyle\sum_{\ell=1}^{2025}\dfrac{a_\ell}{10^\ell}$ with $a_\ell\in\{0,1,\ldots,9\}$ for $1\le\ell\le2025$ and $a_{2025}\neq0$. Let $\{\beta_n\}_{n\ge1}$ be a strictly increasing sequence in $(0,1)$ converging to $\alpha$. For each $n\ge1$, write $\beta_n=\sum_{\ell=1}^\infty\dfrac{b_{n,\ell}}{10^\ell}$ with $b_{n,\ell}\in\{0,1,\ldots,9\}$. Which of the following are necessarily true?
\begin{enumerate}[label=\arabic*.]
\item There exists a positive integer $N$ such that for all $n\ge N$, $b_{n,\ell}=a_\ell$ for all $1\le\ell\le2023$.
\item There exists a positive integer $N$ such that for all $n\ge N$, $b_{n,2025}=a_{2025}-1$.
\item $\beta_n$ is rational for infinitely many $n$.
\item There exists a positive integer $N$ such that for all $n\ge N$, $b_{n,2024}=a_{2024}-1$.
\end{enumerate}\ans{1, 2}

\textbf{Q.67} Let $\mathcal{F}$ be the set of all functions $f:[0,\infty)\to[0,\infty)$ that are Riemann integrable on $[0,t]$ for all $t\in[0,\infty)$. Which of the following are necessarily true?
\begin{enumerate}[label=\arabic*.]
\item If $f\in\mathcal{F}$ is a continuous function and $f(n)=0$ for all positive integers $n$, then $\lim_{x\to\infty}\int_0^xf(t)\,dt$ exists in $\mathbb{R}$.
\item If $f\in\mathcal{F}$ is uniformly continuous and $\lim_{x\to\infty}\int_0^xf(t)\,dt$ exists in $\mathbb{R}$, then $\lim_{x\to\infty}f(x)=0$.
\item There exists a function $f\in\mathcal{F}$ such that $\lim_{x\to\infty}\int_0^xf(t)\,dt$ exists in $\mathbb{R}$ but $\lim_{x\to\infty}f(x)\neq0$.
\item If $f\in\mathcal{F}$ is Lebesgue measurable and $\lim_{x\to\infty}\int_0^xf(t)\,dt=0$, then $f=0$ a.e.\ on $[0,\infty)$.
\end{enumerate}\ans{2, 3, 4}

\textbf{Q.68} Consider the ODE $y''+r(x)y=0$, where $r(x)=\begin{cases}x^3\sin\!\left(\tfrac{1}{x}\right),&x\neq0\\0,&x=0\end{cases}$. Then which of the following are true?
\begin{enumerate}[label=\arabic*.]
\item There exists a non-trivial solution $\phi$ of ODE, and $\alpha,\beta\in\mathbb{R}$, $\alpha<\beta$, such that $\phi$ has infinitely many zeros in $[\alpha,\beta]$.
\item There exists a non-trivial solution $\phi$ of ODE, and a sequence $\{x_n\}$ such that $x_n\to\infty$ and $\phi(x_n)=0$ for all $n\in\mathbb{N}$.
\item If $\phi_1,\phi_2$ are two linearly independent solutions of ODE, then their Wronskian is a constant function on $\mathbb{R}$.
\item There exists a non-trivial solution of ODE which vanishes at most at one point in $(0,\infty)$.
\end{enumerate}\ans{2, 3}

\textbf{Q.69} Consider $X=\prod_{c\in(0,\infty)}[0,c]$ together with the product topology, where $[0,c]$ is equipped with the Euclidean topology. Which of the following are necessarily true?
\begin{enumerate}[label=\arabic*.]
\item $X$ is metrizable. \item $X$ is compact. \item $X$ is Hausdorff. \item $X$ is connected.
\end{enumerate}\ans{2, 3, 4}

\textbf{Q.70} Let $X_1,X_2,\ldots,X_n$ ($n\ge2$) be a random sample from Uniform$\!\bigl[-\tfrac{3\theta}{2},\tfrac{\theta}{2}\bigr]$ distribution, $\theta>0$ unknown. Let $\bar{X}$, $X_{(1)}$ and $X_{(n)}$ denote the sample mean, smallest and largest order statistics, respectively. Which of the following are true?
\begin{enumerate}[label=\arabic*.]
\item The method of moments estimator of $\theta$ is $-2\bar{X}$
\item The method of moments estimator of $\theta$ is $2\bar{X}$
\item The maximum likelihood estimator of $\theta$ is $\min\!\Bigl\{-\dfrac{2X_{(1)}}{3},2X_{(n)}\Bigr\}$
\item The maximum likelihood estimator of $\theta$ is $\max\!\Bigl\{-\dfrac{2X_{(1)}}{3},2X_{(n)}\Bigr\}$
\end{enumerate}\ans{1, 4}

\textbf{Q.71} Let $\tau_E$ denote the Euclidean topology on $\mathbb{R}$. Which of the following are necessarily true?
\begin{enumerate}[label=\arabic*.]
\item $(\mathbb{R},\tau_E)$ is a normal space.
\item Let $\tau$ be a topology on $\mathbb{R}$. If the identity function of $\mathbb{R}$ is a continuous map from $(\mathbb{R},\tau_E)$ to $(\mathbb{R},\tau)$, then $(\mathbb{R},\tau)$ is a regular space.
\item Let $\tau$ be a topology on $\mathbb{R}$. If the identity function of $\mathbb{R}$ is a continuous map from $(\mathbb{R},\tau)$ to $(\mathbb{R},\tau_E)$, then $(\mathbb{R},\tau)$ is a Hausdorff space.
\item $\mathbb{R}$ is a regular space in the finite-complement topology.
\end{enumerate}\ans{1, 3}

\textbf{Q.72} Let $y(x)$ be the solution of the integral equation $y(x)=e^x+e+1+\tfrac{1}{3}\int_0^1y(t)\,dt$, and let $R\!\left(x,t,\tfrac{1}{3}\right)$ be the resolvent kernel associated to the IE. Then which of the following are true?
\begin{enumerate}[label=\arabic*.]
\item $R\!\left(0,1,\tfrac{1}{3}\right)=\dfrac{3}{2}$ \item $R\!\left(\tfrac{1}{2},\tfrac{1}{2},\tfrac{1}{3}\right)=\dfrac{2}{3}$
\item $y(1)=3e+1$ \item $y(1)=e+3$
\end{enumerate}\ans{1, 3}

\textbf{Q.73} Suppose $X|\theta\sim\text{Binomial}(7,\theta)$, $0<\theta<1$, and the prior distribution of $\theta$ is Beta$(\alpha,\beta)$ where $\alpha>0$ and $\beta>0$ are known. Which of the following statements \textbf{may not} be true?
\begin{enumerate}[label=\arabic*.]
\item Posterior mean of $\theta$ given $X=2$ is less than the prior mean
\item Posterior mean of $\theta$ given $X=3$ is greater than the prior mean
\item Posterior mean of $\theta$ given $X=4$ is not equal to the prior mean
\item Posterior mean of $\theta$ given $X=5$ is equal to the prior mean
\end{enumerate}\ans{1, 2, 3, 4}

\textbf{Q.74} Let $A\in M_4(\mathbb{C})$ be such that $A^2=I$ and $\text{Trace}(A)=0$. Which of the following are necessarily true?
\begin{enumerate}[label=\arabic*.]
\item The set $\{B\in M_4(\mathbb{C})\mid AB+BA=0\}$ is an 8-dimensional $\mathbb{C}$-vector space.
\item The set $\{B\in M_4(\mathbb{C})\mid AB-BA=0\}$ is an 8-dimensional $\mathbb{C}$-vector space.
\item There exists $B\in M_4(\mathbb{C})$ such that $AB+BA=0$, $B^2=I$ and $\text{Trace}(B)=0$.
\item If $B\in M_4(\mathbb{C})$ is such that $AB+BA=0$, then $B$ is diagonalisable.
\end{enumerate}\ans{1, 2, 3}

\textbf{Q.75} Consider an $M/M/2$ queuing system with birth rate $\lambda=4$ per minute, death rate $\mu=1$ per minute, and total capacity of 3 customers. Let $p_0$ and $p_3$ denote the long-run probabilities that the system will be empty and full, respectively. Which of the following are true?
\begin{enumerate}[label=\arabic*.]
\item $p_0+p_3=\dfrac{17}{29}$ \item $p_0>p_3$ \item $\dfrac{p_3}{p_0}=2$ \item $p_3-p_0=\dfrac{15}{29}$
\end{enumerate}\ans{1, 4}

\textbf{Q.76} For $n\ge3$, consider the space $M_n(\mathbb{C})$ of $n\times n$ complex matrices endowed with the inner product $\langle A,B\rangle=\text{Trace}(A^*B)$. For $0\le k\le n$, let $W_k=\{A\in M_n(\mathbb{C}):\langle A,B\rangle=0\text{ for all }B\in M_n(\mathbb{C})\text{ with rank }k\}$. Which of the following are necessarily true?
\begin{enumerate}[label=\arabic*.]
\item $W_0=\{0\}$ \item $W_1=\{0\}$ \item $W_2=\{0\}$ \item $W_n=\{0\}$
\end{enumerate}\ans{2, 3, 4}

\textbf{Q.77} There are six persons. On each day of the year 2024 at least one of them borrowed money from another. A person $P$ is labeled \textit{troublesome} in calendar month $M$ of 2024 if $P$ borrowed from the same person at least twice during the month $M$. Which of the following are necessarily true?
\begin{enumerate}[label=\arabic*.]
\item In every calendar month of 2024, at least one person was troublesome.
\item In any two consecutive calendar months of 2024, at least one person was troublesome in at least one of the two calendar months.
\item At least one person was troublesome in two calendar months of 2024.
\item At least one person was troublesome in January 2024.
\end{enumerate}\ans{2, 3, 4}

\textbf{Q.78} Let $V$ be a real vector space. Suppose $\{u,v,w,x,y\}\subseteq V$ is a spanning set of $V$ and that $\{u,v,x,y\}$ is linearly independent. Which of the following are necessarily true?
\begin{enumerate}[label=\arabic*.]
\item The dimension of $V$ is 4 or 5.
\item If $2u-3v+5w=0$, then $\{v,w,x,y\}$ is a basis of $V$.
\item If $u-6v+7w=0$, then the span of $\{v,w,x\}$ is a 2-dimensional vector space.
\item If $u+w=v+x$, then $\{u,v,w,y\}$ is a basis of $V$.
\end{enumerate}\ans{1, 2, 4}

\textbf{Q.79} Suppose $y(x)$ is the extremal of the variational problem $J(y)=\int_0^1x^2(y')^2\,dx$ subject to $y(0)=0$, $y(1)=1$, $\int_0^1y^2\,dx=\tfrac{1}{7}$. Then which of the following are true?
\begin{enumerate}[label=\arabic*.]
\item $y'\!\left(\tfrac{1}{2}\right)=\dfrac{3}{4}$
\item $y'\!\left(\tfrac{1}{\sqrt{3}}\right)=1$
\item $y'\!\left(\tfrac{1}{3}\right)=\dfrac{1}{3}$
\item $y'\!\left(\tfrac{1}{4}\right)=\dfrac{1}{2}$
\end{enumerate}\ans{1, 2, 3}

\textbf{Q.80} Let $X_1,X_2,\ldots,X_{20}$ be a random sample from $N_{12}(\mu,\Sigma)$, where $\mu\in\mathbb{R}^{12}$ and $\Sigma$ is positive definite. Suppose $\bar{X}=\tfrac{1}{20}\sum_{i=1}^{20}X_i$ and $S=\tfrac{1}{19}\sum_{i=1}^{20}(X_i-\bar{X})(X_i-\bar{X})^T$. Then which of the following are true?
\begin{enumerate}[label=\arabic*.]
\item $19(\bar{X}^TS\bar{X})(\bar{X}^T\Sigma\bar{X})^{-1}\sim\chi^2_{19}$
\item $E(S^{-1})=\dfrac{19}{6}\Sigma^{-1}$
\item $S\sim W_{12}(19,19\Sigma)$
\item $\text{trace}(19\Sigma^{-1}S)\sim\chi^2_{228}$
\end{enumerate}\ans{1, 2, 4}

\textbf{Q.81} Let $u(x,t)$ be the solution of the initial-boundary value problem $u_{tt}-u_{xx}+u_t=0$, $0<x<\pi$, $t>0$; $u(0,t)=0$, $u(\pi,t)=0$, $t>0$; $u(x,0)=\sin x$, $u_t(x,0)=0$, $0<x<\pi$. For $t\ge0$, define $E(t)=\int_0^\pi(u_t^2+u_x^2)\,dx$. Then which of the following are true?
\begin{enumerate}[label=\arabic*.]
\item $E(t)\ge\pi$ for all $t>0$. \item $E(t)\le\pi$ for all $t>0$.
\item $E(t)\ge\dfrac{\pi}{2}$ for all $t>0$. \item $E(t)\le\dfrac{\pi}{2}$ for all $t>0$.
\end{enumerate}\ans{2, 4}

\textbf{Q.82} Let $S$ be the set of all $2\times2$ matrices $A$ such that the iterative sequence generated by the Gauss-Seidel method converges for every initial guess when employed to solve the system $A\begin{pmatrix}x_1\\x_2\end{pmatrix}=\begin{pmatrix}1\\2\end{pmatrix}$. Then which of the following are true?
\begin{enumerate}[label=\arabic*.]
\item $\begin{pmatrix}4&1\\2&3\end{pmatrix}\in S$
\item $\begin{pmatrix}1&2\\3&4\end{pmatrix}\in S$
\item $\begin{pmatrix}1&5\\1&10\end{pmatrix}\in S$
\item $\begin{pmatrix}-5&2\\1&-4\end{pmatrix}\in S$
\end{enumerate}\ans{1, 3, 4}

\textbf{Q.83} Consider a commutative ring $R$ with unity with at least five elements such that for any two elements $a,b\in R$, there exists $c\in R$ such that $a=bc$ or $b=ac$. Which of the following are necessarily true?
\begin{enumerate}[label=\arabic*.]
\item $R$ has a unique maximal ideal. \item Every finitely generated ideal of $R$ is principal.
\item $R$ is an integral domain. \item $R$ is a Euclidean domain.
\end{enumerate}\ans{1, 2}

\textbf{Q.84} Consider the following statements:\\
\textbf{(P)} The IVP $y'+y=e^{-x^2}$, $y(0)=0$, has a Taylor series solution about $x=0$.\\
\textbf{(Q)} The IVP $y'+y=r(x)$, $y(0)=0$, where $r(x)=\begin{cases}e^{-1/x^2},&x\neq0\\0,&x=0\end{cases}$, has a Taylor series solution about $x=0$.\\
Then which of the following are true?
\begin{enumerate}[label=\arabic*.]
\item (P) is true. \item (P) is FALSE.
\item (Q) is true. \item (Q) is FALSE.
\end{enumerate}\ans{1, 4}

\textbf{Q.85} Let $F:\mathbb{R}^3\to\mathbb{R}$ be continuously differentiable such that $F(0,0,0)=-1$, $\dfrac{\partial F}{\partial x}(0,0,0)=2$, $\dfrac{\partial F}{\partial y}(0,0,0)=0$, and $\dfrac{\partial F}{\partial z}(0,0,0)=3$. For $\epsilon>0$, define $\Omega_\epsilon=(-\epsilon,\epsilon)\times(-\epsilon,\epsilon)\subset\mathbb{R}^2$. Which of the following are necessarily true?
\begin{enumerate}[label=\arabic*.]
\item There exist $\epsilon,\delta>0$ and continuously differentiable $g:\Omega_\epsilon\to(-\delta,\delta)$ such that $g(0,0)=0$ and $F(x,g(x,y),z)=-1$ for all $(x,y)\in\Omega_\epsilon$.
\item There exist $\epsilon,\delta>0$ and continuously differentiable $h:\Omega_\epsilon\to(-\delta,\delta)$ such that $h(0,0)=0$ and $F(x,h(x,z),z)=-1$ for all $(x,z)\in\Omega_\epsilon$.
\item There exist $\epsilon,\delta>0$ and continuously differentiable functions $k_1,k_2:(-\epsilon,\epsilon)\to(-\delta,\delta)$ such that $k_1(0)=k_2(0)=0$ and $F(x,k_1(x),k_2(x))=-1$ for all $x\in(-\epsilon,\epsilon)$.
\item There exist $\epsilon,\delta>0$ and continuously differentiable functions $j_1,j_2:(-\epsilon,\epsilon)\to(-\delta,\delta)$ such that $j_1(0)=j_2(0)=0$ and $F(j_1(z),j_2(z),z)=-1$ for all $z\in(-\epsilon,\epsilon)$.
\end{enumerate}\ans{1, 3, 4}

\textbf{Q.86} For each positive integer $n$, let $f_n:[-1,1]\to\mathbb{R}$ be given by $f_n(x)=\begin{cases}-1,&-1\le x\le-\tfrac{1}{n}\\nx,&-\tfrac{1}{n}<x<\tfrac{1}{n}\\1,&\tfrac{1}{n}\le x\le1\end{cases}$. On which of the following intervals does the sequence $\{f_n\}_{n\ge1}$ converge uniformly?
\begin{enumerate}[label=\arabic*.]
\item $[0,1]$ \item $(0,1]$ \item $(10^{-2025},1)$ \item $(-10^{-2025},10^{-2025}]$
\end{enumerate}\ans{3, 4}

\textbf{Q.87} Let $(X,d)$ be a metric space. For a non-empty subset $A$ of $X$, and $x\in X$, define $d(x,A)=\inf_{a\in A}d(x,a)$. Which of the following are necessarily true?
\begin{enumerate}[label=\arabic*.]
\item For all $x,y\in X$ and every non-empty subset $A$ of $X$, $d(x,A)-d(y,A)\le d(x,y)$.
\item For every non-empty subset $A$ of $X$, the function $x\mapsto d(x,A)$ is uniformly continuous on $X$.
\item A non-empty subset $A$ of $X$ is closed if and only if $d(x,A)>0$ for all $x\in X\setminus A$.
\item If $X$ is compact and $C_1,\ldots,C_k$ are non-empty closed sets of $X$ such that $\bigcap_{1\le i\le k}C_i=\emptyset$, then the minimum value of the function $x\mapsto\sum_{1\le i\le k}d(x,C_i)$ on $X$ is 0.
\end{enumerate}\ans{1, 2, 3}

\textbf{Q.88} Consider the following subset of real numbers: $A=\bigl\{(1+(-1)^n)n-\tfrac{1}{n}:n\text{ is a positive integer}\bigr\}$. Which of the following are true?
\begin{enumerate}[label=\arabic*.]
\item $A$ is bounded below but not bounded above.
\item $A$ is bounded above but not bounded below.
\item $\inf A=-1$ \item $\sup A=1$
\end{enumerate}\ans{1, 3}

\textbf{Q.89} Let $f:(0,\infty)\to(0,\infty)$ be the function defined by $f(x)=\int_0^x\dfrac{\sqrt{t}}{1+t^2}\,dt$, where $\sqrt{t}$ denotes the positive square root for $t>0$. Which of the following are true?
\begin{enumerate}[label=\arabic*.]
\item $f$ is a uniformly continuous function. \item $f$ is a bounded function.
\item There exists $x\in(0,\infty)$ such that $f(x)=0$. \item The derivative of $f$ is continuous.
\end{enumerate}\ans{1, 2, 4}

\textbf{Q.90} Consider the following real-valued functions $F_1$ and $F_2$ defined on $\mathbb{R}$, given by $F_1(x)=\begin{cases}1,&x\ge1\\0,&\text{otherwise}\end{cases}$, $F_2(x)=\begin{cases}\int_0^xe^{-t}\,dt,&x\ge0\\0,&\text{otherwise}\end{cases}$. Define $F:\mathbb{R}\to[0,1]$ by $F(x)=\tfrac{2}{3}F_1(x)+\tfrac{1}{3}F_2(x)$ for all $x\in\mathbb{R}$. Which of the following are true?
\begin{enumerate}[label=\arabic*.]
\item $F$ is non-decreasing on $\mathbb{R}$ \item $\lim_{x\to\infty}F(x)=1$
\item $F$ is left-continuous on $\mathbb{R}$ \item $F$ is right-continuous on $\mathbb{R}$
\end{enumerate}\ans{1, 2, 4}

\textbf{Q.91} Consider the following $2^5$-factorial design with 8 blocks. [\textit{Block structure as given in original question}] Which of the following are true?
\begin{enumerate}[label=\arabic*.]
\item $ABC$ is confounded with blocks \item $BCD$ is confounded with blocks
\item $CDE$ is confounded with blocks \item $ABDE$ is confounded with blocks
\end{enumerate}\ans{1, 2, 3, 4}

\textbf{Q.92} Let $\mathbb{N}$ denote the set of all positive integers. For $n\in\mathbb{N}$, let $f_n:[0,1]\to\mathbb{R}$ be given by $f_n(x)=\begin{cases}n(1-nx),&0\le x\le\tfrac{1}{n}\\0,&x>\tfrac{1}{n}\end{cases}$. Which of the following are necessarily true?
\begin{enumerate}[label=\arabic*.]
\item The sequence $\{f_n\}_{n\ge1}$ is uniformly bounded.
\item The sequence $\{f_n\}_{n\ge1}$ does not converge pointwise.
\item The sequence $\{f_n\}_{n\ge1}$ converges uniformly to the constant function 0.
\item The set $\{f_n:n\in\mathbb{N}\}$ is compact in $C[0,1]$ equipped with the supremum norm.
\end{enumerate}\ans{1, 3, 4}

\textbf{Q.93} Let $f$ be an entire function. Consider the function $g$ given by $g(z)=f(z)-\tfrac{1}{z}$ for $z\in\mathbb{C}\setminus\{0\}$. Which of the following are necessarily true?
\begin{enumerate}[label=\arabic*.]
\item The function $g$ has a pole at 0.
\item If $g(\alpha)=0$, then $|\alpha|\neq1$.
\item The function $g$ has only finitely many zeros.
\item $\max_{|z|=1}|g(z)|\ge1$
\end{enumerate}\ans{1, 2, 3, 4}

\textbf{Q.94} Let $\ell^2$ denote the vector space of all square summable sequences $\{a_n\}_{n\ge1}$ of real numbers with the inner product $\langle\{a_n\},\{b_n\}\rangle=\sum_{n=1}^\infty a_nb_n$. Let $H=\bigl\{\{a_n\}\in\ell^2:|a_n|\le\tfrac{1}{n}\text{ for all positive integers }n\bigr\}$. Which of the following are true?
\begin{enumerate}[label=\arabic*.]
\item $H$ contains an orthonormal basis of $\ell^2$. \item $H$ is a linear subspace of $\ell^2$.
\item $H$ is a bounded subset of $\ell^2$. \item $H$ is a convex subset of $\ell^2$.
\end{enumerate}\ans{3, 4}

\textbf{Q.95} Let $\{X_n:n\ge1\}$ be a sequence of i.i.d.\ random variables, where $X_1$ has an Exponential distribution with mean 1. Define $T_n=\max\{X_1,X_2,\ldots,X_n\}-\ln n$, $n\ge1$. Suppose $T_n\xrightarrow{d}Y$ as $n\to\infty$. Then which of the following are true?
\begin{enumerate}[label=\arabic*.]
\item $Y$ has a Double Exponential distribution with location parameter $\ln(\ln2)$ and scale parameter 1
\item Median of $Y=\ln(\ln2)$
\item $P(Y\le0)=e^{-1}$
\item The derivative of the CDF of $Y$ at $\ln3$ is $e^{-1/3}$
\end{enumerate}\ans{3}

\textbf{Q.96} Let $f:[0,1]\to[0,1]$ be a monotonically increasing function. For any $\alpha\in(0,1)$, let $L_\alpha^+=\lim_{x\to\alpha^+}f(x)$ and $L_\alpha^-=\lim_{x\to\alpha^-}f(x)$ denote the right and left hand limits respectively, provided they exist. For $\alpha\in(0,1)$, if $L_\alpha^+$ and $L_\alpha^-$ exist, define $U_\alpha=\begin{cases}(L_\alpha^-,L_\alpha^+)&\text{if }L_\alpha^-<L_\alpha^+\\\emptyset&\text{if }L_\alpha^+\le L_\alpha^-\end{cases}$. Which of the following are true?
\begin{enumerate}[label=\arabic*.]
\item $L_\alpha^+$ and $L_\alpha^-$ exist for every $\alpha\in(0,1)$.
\item If $f$ is surjective, then $f$ is continuous.
\item If $f$ is Riemann integrable, then $f$ is continuous.
\item If the left and right hand limits exist at $\alpha,\beta\in(0,1)$, $\alpha\neq\beta$, then $U_\alpha\cap U_\beta=\emptyset$.
\end{enumerate}\ans{1, 2, 4}

\textbf{Q.97} Suppose that the transition probability matrix of a homogeneous Markov chain with state space $\{1,2,3,4\}$ is given by $P=\begin{pmatrix}\tfrac{1}{4}&\tfrac{3}{4}&0&0\\1&0&0&0\\\tfrac{1}{8}&0&\tfrac{7}{8}&0\\0&0&\tfrac{1}{9}&\tfrac{8}{9}\end{pmatrix}$. Which of the following are true?
\begin{enumerate}[label=\arabic*.]
\item State 2 is a positive recurrent state \item Mean recurrence time of state 1 is $\tfrac{7}{4}$
\item State 4 is a transient state \item State 3 is aperiodic and ergodic
\end{enumerate}\ans{1, 2, 3}

\textbf{Q.98} Let $X$ and $Y$ be two independent random variables such that the moment generating functions of $X$ and $Y$ are $M_X(t)=e^{3(e^t-1)}$, $t\in\mathbb{R}$, and $M_Y(t)=\bigl(\tfrac{1}{2}e^{-3t}+\tfrac{1}{2}e^{3t}\bigr)^2$, $t\in\mathbb{R}$, respectively. Then which of the following are true?
\begin{enumerate}[label=\arabic*.]
\item $P(XY=0)=\dfrac{1+e^{-3}}{2}$ \item $E(X+Y)=3$
\item $\text{Var}(X+Y)=21$ \item $\text{Cov}(X+Y,X-Y)=0$
\end{enumerate}\ans{1, 2, 3}

\textbf{Q.99} Let $A$ be a non-zero $3\times3$ matrix with integer entries. Let $\lambda_i\in\mathbb{C}$, $1\le i\le3$, be all the eigenvalues of $A$ (not necessarily distinct). Which of the following are necessarily true?
\begin{enumerate}[label=\arabic*.]
\item There exists a cubic polynomial $f(X)\in\mathbb{Q}[X]$ such that $f(\lambda_i)=0$ for all $1\le i\le3$.
\item There exists a quadratic polynomial $f(X)\in\mathbb{Q}[X]$ such that $f(\lambda_i)=0$ for all $1\le i\le3$.
\item If $f(X)\in\mathbb{Q}[X]$ is such that $f(\lambda_i)=0$ for all $1\le i\le3$, then $f(A)=0$.
\item If $f(X)\in\mathbb{Q}[X]$ is a cubic polynomial such that $f(\lambda_i)=0$ for all $1\le i\le3$, then $f(A)=0$.
\end{enumerate}\ans{3, 4}

\textbf{Q.100} A mechanical system is described using generalized position $q$ and generalized momentum $p$. Let $Q$ and $P$ denote new generalized position and generalized momentum variables respectively, generated by the generating function $F(q,P)=q^2e^P$, and $Q$, $P$ are canonical coordinates. Let $G(p,Q)$ be a function such that $G(2,e)=0$, and it generates the same canonical coordinates $Q$, $P$. Then which of the following are true?
\begin{enumerate}[label=\arabic*.]
\item $G(p,Q)=-Q\!\left[1+\log_e\!\left(\dfrac{p^2}{4Q}\right)\right]$
\item $G(p,Q)=-pQ\!\left[1+\log_e\!\left(\dfrac{Q^2}{4p}\right)\right]$
\item $p=2qe^P$, $Q=q^2e^P$ \item $p=-2qe^P$, $Q=-q^2e^P$
\end{enumerate}\ans{1, 3}

\textbf{Q.101} Let $\mathbb{D}=\{z\in\mathbb{C}:|z|<1\}$ and $f:\mathbb{D}\to\mathbb{D}$ be a holomorphic function which satisfies $f(-1/2)=0$. Which of the following are necessarily true?
\begin{enumerate}[label=\arabic*.]
\item $|f(-1/5)|\le1/5$ \item $|f(-1/5)|\le1/3$
\item $|f'(-1/2)|\le1/2$ \item $|f'(-1/2)|\le4/3$
\end{enumerate}\ans{2, 4}

\textbf{Q.102} For $t\in[0,2\pi]$, let $\gamma_1(t)=e^{it}$ and $\gamma_2(t)=1+i+e^{it}$. Which of the following are true?
\begin{enumerate}[label=\arabic*.]
\item $\displaystyle\int_{\gamma_1}\dfrac{\mathrm{d}z}{z\sin z}=0$
\item $\displaystyle\int_{\gamma_1}\dfrac{\mathrm{d}z}{z\sin z}=2\pi i$
\item $\displaystyle\int_{\gamma_2}\dfrac{\mathrm{d}z}{z\sin z}=2\pi i$
\item $\displaystyle\int_{\gamma_2}\dfrac{\mathrm{d}z}{z\sin z}=0$
\end{enumerate}\ans{1, 4}

\textbf{Q.103} Let $V$ be a finite-dimensional $\mathbb{R}$-vector space and $T:V\to V$ a linear operator such that $T^2$ is diagonalizable over $\mathbb{R}$. Which of the following are necessarily true?
\begin{enumerate}[label=\arabic*.]
\item If $T$ is not diagonalizable over $\mathbb{R}$, then $T^2$ has an eigenvalue $\le0$.
\item If $T^2$ has only negative eigenvalues, then $\dim V$ is an even integer.
\item If $T^2$ has only non-negative eigenvalues, then $T$ is diagonalizable over $\mathbb{R}$.
\item For each non-zero $v\in V$, $\{v,Tv,T^2v\}$ is linearly dependent.
\end{enumerate}\ans{1, 2}

\textbf{Q.104} Let $X_1,X_2,\ldots,X_7$ be a random sample drawn from a continuous distribution with unknown unique median $M$. The null hypothesis $H_0:M=2$ is tested against the alternative $H_1:M>2$ at level of significance 0.05 using the right-tailed test based on the Sign test statistic $K$, which is the number of observations in the sample greater than 2. If the observed sample is $-3,-6,1,9,4,10,12$, which of the following are true?
\begin{enumerate}[label=\arabic*.]
\item $H_0$ is rejected
\item The $p$-value of the test is greater than 0.01
\item Under $H_0$, $E(K)=3.5$
\item If $X_{(i)}$ denotes the $i$-th smallest observation of the sample, then $[X_{(2)},X_{(6)}]$ is a confidence interval for $M$ with confidence coefficient at least 0.95
\end{enumerate}\ans{2, 3}

\textbf{Q.105} Suppose $y_1(x)$ and $y_2(x)$ are two linearly independent solutions of the differential equation $x^2y''+\dfrac{(1+\sin x)x}{2}y'-\dfrac{3}{2}(\cos x)y=0$, $x>0$, satisfying $y_2(0)=0$. Then which of the following are true?
\begin{enumerate}[label=\arabic*.]
\item $\lim_{x\to0^+}\dfrac{y_2(x)}{x^2y_1(x)}$ exists.
\item $\lim_{x\to0^+}\dfrac{y_2(x)}{xy_1(x)}$ exists.
\item $\lim_{x\to0^+}\dfrac{xy_1(x)}{y_2(x)}$ does NOT exist.
\item $\lim_{x\to0^+}\dfrac{x^2y_1(x)}{y_2(x)}$ exists.
\end{enumerate}\ans{1, 2, 3}

\textbf{Q.106} Suppose two fair dice are thrown independently at random. Let $X$ and $Y$ be the numbers on the upper face of the first die and that of the second die, respectively. Then which of the following are true?
\begin{enumerate}[label=\arabic*.]
\item $P(X-Y=0\mid X+Y=2)=P(X-Y=0\mid X+Y=12)$
\item $E\!\left(\dfrac{X-Y}{X+Y}\right)=0$
\item $\text{Cov}(X+Y,X-Y)=0$
\item $(X+Y)$ and $(X-Y)$ are independent
\end{enumerate}\ans{1, 2, 3}

\textbf{Q.107} Let $V$ be a 4-dimensional complex vector space and $A$ a linear operator on $V$. Which of the following are necessarily true?
\begin{enumerate}[label=\arabic*.]
\item There exist $\lambda\in\mathbb{C}$ and a non-zero $v\in V$ such that $Av=\lambda v$.
\item There exist $\lambda,\mu\in\mathbb{C}$ and linearly independent vectors $v,w\in V$ such that $Av=\lambda v$ and $Aw=\mu w$.
\item There exist $\lambda,\mu,\delta\in\mathbb{C}$ and linearly independent vectors $v,w\in V$ such that $Av=\lambda v$ and $Aw=\mu v+\delta w$.
\item There exists a three-dimensional subspace $W$ such that $Aw\in W$ for all $w\in W$.
\end{enumerate}\ans{1, 3, 4}

\textbf{Q.108} Let $p\ge3$ be a prime number and $f(x)\in\mathbb{Q}[x]$ an irreducible polynomial of degree $p$. Suppose that $a_1,\ldots,a_p\in\mathbb{C}$ are the roots of $f$ and that $a_1\notin\mathbb{R}$, $a_2\notin\mathbb{R}$, and $a_i\in\mathbb{R}$ for all $3\le i\le p$. Let $K=\mathbb{Q}(a_1,\ldots,a_p)$ be the subfield of $\mathbb{C}$ generated by the roots of $f$. Consider the Galois group $G$ of $K$ over $\mathbb{Q}$ as a subgroup of $S_p$, the group of permutations of $\{a_1,\ldots,a_p\}$. Which of the following are true?
\begin{enumerate}[label=\arabic*.]
\item The transposition $(1\,2)$ belongs to $G$. \item $|G|$ is divisible by $p$.
\item A $p$-cycle belongs to $G$. \item $G=S_p$
\end{enumerate}\ans{1, 2, 3, 4}

\textbf{Q.109} Let $V$ be a real vector space and $L(V)$ denote the space of linear operators on $V$. Let $T\in L(V)$ be a non-zero operator such that $T^2=T$. Consider the subspace $W$ of $L(V)$ spanned by $\{I,T^n:n\text{ is a positive integer}\}$. Which of the following are necessarily true?
\begin{enumerate}[label=\arabic*.]
\item The set $\{S\in W:S^2=S\}$ contains exactly 2 elements.
\item $\dim_\mathbb{R}(W)=2$
\item If $U\in L(V)$ is such that $U^2=U$ and $(T+U)^2=T+U$, then $TU=0$.
\item If $U\in L(V)$ is such that $U^2=U$ and $(T-U)^2=T-U$, then $(TU)^2=TU$.
\end{enumerate}\ans{2, 4}

\textbf{Q.110} Let $G$ be a finite non-abelian group. Which of the following are necessarily true?
\begin{enumerate}[label=\arabic*.]
\item If $d$ is a positive integer that divides $|G|$, then $G$ has a subgroup of order $d$.
\item The map $f:G\times G\to G$ given by $f(a,b)=ab$ is not a group homomorphism.
\item Suppose that for every positive integer $d$ that divides $|G|$, there exists a subgroup of $G$ of order $d$. Then $G$ has at least three normal subgroups.
\item $|G|\neq16$
\end{enumerate}\ans{2, 3}

\textbf{Q.111} Consider a series system comprising four components, having independent and identically distributed lifetimes with hazard rate $\lambda(t)=\dfrac{1}{1+t}$, $t>0$, and survival function $S(t)$, $t>0$. If $Y$ denotes the lifetime of the series system, then which of the following are true?
\begin{enumerate}[label=\arabic*.]
\item Cumulative hazard function of each component is $H(t)=2\ln(1+t)$, $t>0$
\item $S(t)=\dfrac{1}{(1+t)^2}$, $t>0$
\item $P(Y<\tfrac{1}{2})=\{S(\tfrac{1}{2})\}^4$
\item $P(Y<\tfrac{1}{2})=\dfrac{65}{81}$
\end{enumerate}\ans{4}

\textbf{Q.112} Let $a,b$ be distinct positive real numbers. Consider the function $f:\mathbb{R}^2\to\mathbb{R}$ given by $f(x,y)=\begin{cases}\dfrac{(ax+by)^2}{ax^2+by^2}&\text{if }(x,y)\neq(0,0)\\0&\text{if }(x,y)=(0,0)\end{cases}$. Which of the following are necessarily true?
\begin{enumerate}[label=\arabic*.]
\item $\lim_{(x,y)\to(0,0)}f(x,y)$ does not exist.
\item The partial derivatives of $f$ at $(0,0)$ do not exist.
\item $\lim_{x\to0}f(x,0)=\lim_{y\to0}f(0,y)$
\item $f$ is differentiable at $(0,0)$.
\end{enumerate}\ans{1}

\textbf{Q.113} Let $X_1,X_2,\ldots,X_n$ ($n>5$) be independent random variables such that $X_t=\alpha+\alpha^2t+\epsilon_t$, $t=1,\ldots,n$, where $\epsilon_1,\epsilon_2,\ldots,\epsilon_n$ are i.i.d.\ $N(0,\sigma^2)$ random variables, $\alpha\in\mathbb{R}$ and $\sigma>0$ are unknown parameters. Which of the following are true?
\begin{enumerate}[label=\arabic*.]
\item $(X_1,X_2,\ldots,X_n)$ is a sufficient statistic for $(\alpha,\sigma)$
\item $\Bigl(\sum_{t=1}^nX_t,\sum_{t=1}^ntX_t,\sum_{t=1}^nt^2X_t^2\Bigr)$ is a jointly minimal sufficient statistic for $(\alpha,\sigma)$
\item $X_4-X_3-X_2+X_1$ is an ancillary statistic
\item $\dfrac{X_4+X_1-X_2-X_3}{X_5+X_2-X_3-X_4}$ is an ancillary statistic
\end{enumerate}\ans{1, 4}

\textbf{Q.114} Suppose $y(x)$ is the extremal of the variational problem $J(y)=\int_0^1\bigl[(y')^2\sin x+(2\cos x)y\bigr]\,dx$ subject to $y(0)=0$, $y(1)=1$. Then which of the following are true?
\begin{enumerate}[label=\arabic*.]
\item $y\!\left(\tfrac{1}{2}\right)=1$ \item $y'(0)=1$
\item $y\!\left(\tfrac{1}{4}\right)=2$ \item $y'\!\left(\tfrac{1}{2}\right)=1$
\end{enumerate}\ans{2, 4}

\textbf{Q.115} Let $Y=X\beta+\epsilon$ be a multiple linear regression model with $p$ regressors and an intercept, where $\beta=(\beta_0,\beta_1,\ldots,\beta_p)^T$ and the random error $\epsilon\sim N_n(0,\sigma^2I)$, $\sigma>0$ and $n>p+1$. Let $\hat{\beta}$ be the least squares estimator of $\beta$. Let the total sum of squares (corrected), sum of squares due to regression/model and sum of squares due to error/residual, based on $\hat{\beta}$, be denoted by $Y^TAY$, $Y^TBY$ and $Y^TCY$, respectively, so that $Y^TAY=Y^TBY+Y^TCY$. Which of the following are always true?
\begin{enumerate}[label=\arabic*.]
\item $\dfrac{Y^TAY}{\sigma^2}$ follows a central $\chi^2$ distribution with $(n-1)$ degrees of freedom.
\item $\dfrac{Y^TBY}{\sigma^2}$ follows a central $\chi^2$ distribution with $p$ degrees of freedom if $\beta_1=\beta_2=\cdots=\beta_p=0$.
\item $\dfrac{Y^TBY}{Y^TCY}$ follows a central $F$ distribution with $(p,n-p)$ degrees of freedom.
\item $Y^TBY$ and $Y^TCY$ are independently distributed if and only if $\beta_1=\beta_2=\cdots=\beta_p=0$.
\end{enumerate}\ans{2}

\textbf{Q.116} Let $\lambda\in\mathbb{R}$ be such that the integral equation $y(x)=\lambda\int_{-1}^1(5xt^3+4x^2t+3xt)y(t)\,dt$ admits a non-trivial solution $y(x)$ such that $y(1)=\tfrac{5}{2}$. Then which of the following are true?
\begin{enumerate}[label=\arabic*.]
\item $y(0)+y'(0)=\dfrac{3}{2}$
\item $y\!\left(\tfrac{1}{2}\right)+y'\!\left(\tfrac{1}{2}\right)=\dfrac{7}{2}$
\item $y(-1)+y'(-1)=-1$
\item $y\!\left(\tfrac{1}{3}\right)+y'\!\left(\tfrac{1}{3}\right)=\dfrac{14}{9}$
\end{enumerate}\ans{1, 2, 3}

\textbf{Q.117} Let $\{x_n\}$ be a convergent iterative sequence generated by Newton-Raphson method for solving $\sin x-1=0$ such that $x_n\to\dfrac{\pi}{2}$ as $n\to\infty$. For $n\in\mathbb{N}$, let $e_n=x_n-\dfrac{\pi}{2}$. Let $p>0$ be such that $\lim_{n\to\infty}\dfrac{|e_{n+1}|}{|e_n|^p}$ exists and is non-zero. Then which of the following are true?
\begin{enumerate}[label=\arabic*.]
\item $p=1$
\item $\lim_{n\to\infty}\dfrac{|e_{n+1}|}{|e_n|^p}=\dfrac{1}{2}$
\item $p=2$
\item $\lim_{n\to\infty}\dfrac{|e_{n+1}|}{|e_n|^p}=1$
\end{enumerate}\ans{1, 2}

\textbf{Q.118} Consider the boundary value problem (BVP) $u_{xx}+u_{yy}=0$ in $\Omega=\{(x,y)\in\mathbb{R}^2:x^2+y^2<1\}$, $u(x,y)=e^{x+y}$ on $\partial\Omega=\{(x,y)\in\mathbb{R}^2:x^2+y^2=1\}$. Then which of the following are true?
\begin{enumerate}[label=\arabic*.]
\item There exists a unique solution to BVP.
\item The BVP does NOT have a solution.
\item There exists a solution $u$ to BVP such that $u(x,y)=\dfrac{1+e}{2}$ for some $(x,y)\in\Omega\cup\partial\Omega$.
\item There exists a solution $u$ to BVP such that $u(x,y)=\dfrac{1+e^3}{2}$ for some $(x,y)\in\Omega\cup\partial\Omega$.
\end{enumerate}\ans{1, 3}

\textbf{Q.119} Consider the simple linear regression model $Y_i=\beta x_i+\epsilon_i$, $i=1,2,\ldots,n$, where $\beta>0$, $\sum_{i=1}^nx_i^2>0$, and the uncorrelated errors $\epsilon_i$ have zero mean and finite variance $\sigma^2(>0)$. Let $\tilde{\beta}_1=\sum_{i=1}^na_i^*Y_i$ minimise $E\!\bigl(\sum_{i=1}^na_iY_i-\beta\bigr)^2$ with respect to scalars $a_1,a_2,\ldots,a_n$. Let $\tilde{\beta}_2$ be the ordinary least squares estimator of $\beta$. Which of the following are true?
\begin{enumerate}[label=\arabic*.]
\item $\text{Var}(\tilde{\beta}_1)>\text{Var}(\tilde{\beta}_2)$, $E(\tilde{\beta}_1-\beta)^2>E(\tilde{\beta}_2-\beta)^2$
\item $\text{Var}(\tilde{\beta}_1)<\text{Var}(\tilde{\beta}_2)$, $E(\tilde{\beta}_1-\beta)^2<E(\tilde{\beta}_2-\beta)^2$
\item $\text{Var}(\tilde{\beta}_1)>\text{Var}(\tilde{\beta}_2)$, $E(\tilde{\beta}_1)<E(\tilde{\beta}_2)$
\item $\text{Var}(\tilde{\beta}_1)<\text{Var}(\tilde{\beta}_2)$, $E(\tilde{\beta}_1)>E(\tilde{\beta}_2)$
\end{enumerate}\ans{2}

\textbf{Q.120} Let $f:\mathbb{C}\to\mathbb{C}$ be defined by $f(z)=\sin^2z+\cos^2|z|$. Which of the following are true?
\begin{enumerate}[label=\arabic*.]
\item $f$ is a real valued function.
\item $f(z)=1$ for all $z\in\mathbb{C}$.
\item $f$ is not an entire function.
\item $f$ has finitely many zeros on the imaginary axis.
\end{enumerate}\ans{3, 4}

\vspace{1em}\hrule
\begin{center}
\textbf{END OF CSIR-UGC NET DECEMBER 2025 --- MATHEMATICAL SCIENCES}
\end{center}

\end{document}
